\chapter{Code Snippets}
\label{appendix:code}

This appendix provides key code snippets from the implementation of the Physics-Informed Neural Network for Energy-efficient Control (PIEC) system. All code is written in Python using PyTorch for neural network components.

\section{PINN Model Forward Pass}
\label{sec:pinn_forward}

The following code implements the forward pass of the Physics-Informed Neural Network, including the computation of physics-informed loss terms based on the equations of motion.

\begin{lstlisting}[language=Python, caption={PINN Model Forward Pass Implementation}, label={lst:pinn_forward}, numbers=left, numberstyle=\tiny, stepnumber=1, basicstyle=\ttfamily\footnotesize, breaklines=true, frame=single]
import torch
import torch.nn as nn
import torch.nn.functional as F

class PINNModel(nn.Module):
    """
    Physics-Informed Neural Network for mobile robot dynamics.
    Incorporates kinematic and dynamic constraints.
    """
    def __init__(self, input_dim=6, hidden_dims=[128, 128, 64], 
                 output_dim=3, dropout_rate=0.1):
        super(PINNModel, self).__init__()
        
        # Build network layers
        layers = []
        prev_dim = input_dim
        for hidden_dim in hidden_dims:
            layers.append(nn.Linear(prev_dim, hidden_dim))
            layers.append(nn.Tanh())
            layers.append(nn.Dropout(dropout_rate))
            prev_dim = hidden_dim
        
        self.hidden_layers = nn.ModuleList(layers)
        self.output_layer = nn.Linear(prev_dim, output_dim)
        
        # Physics parameters
        self.mass = 12.5  # kg
        self.inertia = 0.8  # kg*m^2
        self.wheel_radius = 0.1  # m
        self.wheel_base = 0.3  # m
        
    def forward(self, x, compute_physics_loss=False):
        """
        Forward pass through the network.
        
        Args:
            x: Input tensor [batch, input_dim]
               Format: [x, y, theta, v, omega, energy]
            compute_physics_loss: Whether to compute physics terms
            
        Returns:
            output: Predicted controls [v_cmd, omega_cmd, power]
            physics_loss: Physics-informed loss (if requested)
        """
        # Extract state variables
        pos_x = x[:, 0:1]
        pos_y = x[:, 1:2]
        theta = x[:, 2:3]
        v = x[:, 3:4]
        omega = x[:, 4:5]
        energy = x[:, 5:6]
        
        # Forward through hidden layers
        h = x
        for layer in self.hidden_layers:
            h = layer(h)
        
        # Output layer
        output = self.output_layer(h)
        v_cmd = output[:, 0:1]
        omega_cmd = output[:, 1:2]
        power_pred = output[:, 2:3]
        
        if not compute_physics_loss:
            return output
        
        # Compute physics-informed loss terms
        # Enable gradient computation for physics constraints
        x_phys = x.clone().requires_grad_(True)
        
        # Recompute forward pass for gradient computation
        h_phys = x_phys
        for layer in self.hidden_layers:
            h_phys = layer(h_phys)
        output_phys = self.output_layer(h_phys)
        
        # Extract predicted controls
        v_cmd_phys = output_phys[:, 0:1]
        omega_cmd_phys = output_phys[:, 1:2]
        power_pred_phys = output_phys[:, 2:3]
        
        # Compute state derivatives
        dx_dt = torch.autograd.grad(
            outputs=output_phys.sum(),
            inputs=x_phys,
            create_graph=True,
            retain_graph=True
        )[0]
        
        # Kinematic constraints (differential drive)
        v_x = v * torch.cos(theta)
        v_y = v * torch.sin(theta)
        
        kinematic_loss_x = dx_dt[:, 0:1] - v_x
        kinematic_loss_y = dx_dt[:, 1:2] - v_y
        kinematic_loss_theta = dx_dt[:, 2:3] - omega
        
        # Dynamic constraints (Newton's laws)
        F_linear = self.mass * (v_cmd_phys - v) / 0.1  # dt = 0.1s
        tau_angular = self.inertia * (omega_cmd_phys - omega) / 0.1
        
        dynamic_loss_v = torch.abs(F_linear) - \
                        torch.abs(self.mass * dx_dt[:, 3:4])
        dynamic_loss_omega = torch.abs(tau_angular) - \
                            torch.abs(self.inertia * dx_dt[:, 4:5])
        
        # Energy conservation constraint
        # Power = F*v + tau*omega
        power_true = (F_linear * v_cmd_phys + 
                     tau_angular * omega_cmd_phys)
        energy_loss = torch.abs(power_pred_phys - power_true)
        
        # Combine physics losses
        physics_loss = (
            torch.mean(kinematic_loss_x ** 2) +
            torch.mean(kinematic_loss_y ** 2) +
            torch.mean(kinematic_loss_theta ** 2) +
            torch.mean(dynamic_loss_v ** 2) +
            torch.mean(dynamic_loss_omega ** 2) +
            torch.mean(energy_loss ** 2)
        )
        
        return output, physics_loss
    
    def compute_energy_cost(self, v_cmd, omega_cmd, dt=0.1):
        """
        Compute energy cost for given control commands.
        
        Args:
            v_cmd: Linear velocity command
            omega_cmd: Angular velocity command
            dt: Time step
            
        Returns:
            energy_cost: Estimated energy consumption
        """
        # Motor model: P = k1*v^2 + k2*omega^2 + k3*|v*omega|
        k1 = 0.5  # Linear motion coefficient
        k2 = 0.3  # Rotational motion coefficient
        k3 = 0.2  # Coupling coefficient
        
        power = (k1 * v_cmd**2 + 
                k2 * omega_cmd**2 + 
                k3 * torch.abs(v_cmd * omega_cmd))
        energy_cost = power * dt
        
        return energy_cost
\end{lstlisting}

\section{NSGA-II Non-Dominated Sorting}
\label{sec:nsga2_sorting}

This implementation of the NSGA-II non-dominated sorting algorithm is used for multi-objective path optimization, balancing path length, smoothness, and energy consumption.

\begin{lstlisting}[language=Python, caption={NSGA-II Non-Dominated Sorting Algorithm}, label={lst:nsga2_sorting}, numbers=left, numberstyle=\tiny, stepnumber=1, basicstyle=\ttfamily\footnotesize, breaklines=true, frame=single]
import numpy as np
from typing import List, Tuple, Dict

class NSGA2Optimizer:
    """
    Non-dominated Sorting Genetic Algorithm II (NSGA-II)
    for multi-objective path optimization.
    """
    def __init__(self, population_size=100, max_generations=200,
                 crossover_rate=0.9, mutation_rate=0.1,
                 tournament_size=2):
        self.population_size = population_size
        self.max_generations = max_generations
        self.crossover_rate = crossover_rate
        self.mutation_rate = mutation_rate
        self.tournament_size = tournament_size
        
    def fast_non_dominated_sort(self, objectives):
        """
        Fast non-dominated sorting algorithm.
        
        Args:
            objectives: Array of shape [pop_size, n_objectives]
                       Each row is the objective values for one solution
                       
        Returns:
            fronts: List of lists, each containing indices of 
                   solutions in that Pareto front
        """
        pop_size = len(objectives)
        
        # Initialize domination data structures
        domination_count = np.zeros(pop_size, dtype=int)
        dominated_solutions = [[] for _ in range(pop_size)]
        
        # Find domination relationships
        for i in range(pop_size):
            for j in range(i + 1, pop_size):
                # Check if i dominates j
                i_dominates_j = self._dominates(
                    objectives[i], objectives[j]
                )
                # Check if j dominates i
                j_dominates_i = self._dominates(
                    objectives[j], objectives[i]
                )
                
                if i_dominates_j:
                    dominated_solutions[i].append(j)
                    domination_count[j] += 1
                elif j_dominates_i:
                    dominated_solutions[j].append(i)
                    domination_count[i] += 1
        
        # Build fronts
        fronts = [[]]
        for i in range(pop_size):
            if domination_count[i] == 0:
                fronts[0].append(i)
        
        # Generate subsequent fronts
        current_front = 0
        while len(fronts[current_front]) > 0:
            next_front = []
            for i in fronts[current_front]:
                for j in dominated_solutions[i]:
                    domination_count[j] -= 1
                    if domination_count[j] == 0:
                        next_front.append(j)
            
            current_front += 1
            fronts.append(next_front)
        
        # Remove empty last front
        if len(fronts[-1]) == 0:
            fronts.pop()
        
        return fronts
    
    def _dominates(self, obj1, obj2):
        """
        Check if obj1 dominates obj2 (minimization).
        
        Args:
            obj1: Objective values for solution 1
            obj2: Objective values for solution 2
            
        Returns:
            True if obj1 dominates obj2
        """
        # obj1 dominates obj2 if:
        # - obj1 is no worse in all objectives
        # - obj1 is strictly better in at least one objective
        better_in_all = np.all(obj1 <= obj2)
        better_in_one = np.any(obj1 < obj2)
        
        return better_in_all and better_in_one
    
    def calculate_crowding_distance(self, objectives, front_indices):
        """
        Calculate crowding distance for solutions in a front.
        
        Args:
            objectives: Array of shape [pop_size, n_objectives]
            front_indices: Indices of solutions in this front
            
        Returns:
            distances: Crowding distance for each solution in front
        """
        n_solutions = len(front_indices)
        n_objectives = objectives.shape[1]
        
        # Initialize distances to 0
        distances = np.zeros(n_solutions)
        
        # Calculate crowding distance for each objective
        for obj_idx in range(n_objectives):
            # Sort solutions by this objective
            obj_values = objectives[front_indices, obj_idx]
            sorted_indices = np.argsort(obj_values)
            
            # Boundary solutions get infinite distance
            distances[sorted_indices[0]] = np.inf
            distances[sorted_indices[-1]] = np.inf
            
            # Calculate distance for interior solutions
            obj_range = obj_values[sorted_indices[-1]] - \
                       obj_values[sorted_indices[0]]
            
            if obj_range > 0:
                for i in range(1, n_solutions - 1):
                    distances[sorted_indices[i]] += (
                        (obj_values[sorted_indices[i + 1]] -
                         obj_values[sorted_indices[i - 1]]) /
                        obj_range
                    )
        
        return distances
    
    def tournament_selection(self, population, objectives, 
                           fronts, crowding_distances):
        """
        Select parents using binary tournament selection.
        
        Args:
            population: List of solution chromosomes
            objectives: Objective values for all solutions
            fronts: Pareto fronts from non-dominated sorting
            crowding_distances: Crowding distances for all solutions
            
        Returns:
            selected: List of selected parent solutions
        """
        selected = []
        pop_size = len(population)
        
        # Create front rank mapping
        front_rank = np.zeros(pop_size, dtype=int)
        for rank, front in enumerate(fronts):
            for idx in front:
                front_rank[idx] = rank
        
        # Perform tournament selection
        for _ in range(pop_size):
            # Randomly select tournament candidates
            candidates = np.random.choice(
                pop_size, 
                size=self.tournament_size, 
                replace=False
            )
            
            # Select best candidate based on:
            # 1. Lower front rank (better)
            # 2. Higher crowding distance (more diverse)
            best_idx = candidates[0]
            for candidate_idx in candidates[1:]:
                if (front_rank[candidate_idx] < front_rank[best_idx] or
                    (front_rank[candidate_idx] == front_rank[best_idx] and
                     crowding_distances[candidate_idx] > 
                     crowding_distances[best_idx])):
                    best_idx = candidate_idx
            
            selected.append(population[best_idx].copy())
        
        return selected
    
    def optimize(self, objective_function, bounds, 
                 n_objectives=3, verbose=True):
        """
        Run NSGA-II optimization.
        
        Args:
            objective_function: Function that takes a solution
                              and returns objective values
            bounds: List of (min, max) tuples for each variable
            n_objectives: Number of objectives
            verbose: Whether to print progress
            
        Returns:
            pareto_front: Final Pareto front solutions
            pareto_objectives: Objective values for Pareto front
        """
        n_vars = len(bounds)
        
        # Initialize population randomly
        population = []
        for _ in range(self.population_size):
            individual = np.array([
                np.random.uniform(bounds[i][0], bounds[i][1])
                for i in range(n_vars)
            ])
            population.append(individual)
        
        # Evolution loop
        for generation in range(self.max_generations):
            # Evaluate objectives
            objectives = np.array([
                objective_function(ind) 
                for ind in population
            ])
            
            # Non-dominated sorting
            fronts = self.fast_non_dominated_sort(objectives)
            
            # Calculate crowding distances for all fronts
            all_distances = np.zeros(self.population_size)
            for front in fronts:
                if len(front) > 0:
                    distances = self.calculate_crowding_distance(
                        objectives, front
                    )
                    all_distances[front] = distances
            
            if verbose and generation % 20 == 0:
                print(f"Generation {generation}: "
                      f"Pareto front size = {len(fronts[0])}")
            
            # Selection
            parents = self.tournament_selection(
                population, objectives, fronts, all_distances
            )
            
            # Crossover and mutation
            offspring = self._create_offspring(parents, bounds)
            
            # Combine parent and offspring populations
            combined_pop = population + offspring
            combined_obj = np.vstack([
                objectives,
                np.array([objective_function(ind) 
                         for ind in offspring])
            ])
            
            # Select next generation
            population, objectives = self._select_next_generation(
                combined_pop, combined_obj
            )
        
        # Return final Pareto front
        fronts = self.fast_non_dominated_sort(objectives)
        pareto_indices = fronts[0]
        pareto_front = [population[i] for i in pareto_indices]
        pareto_objectives = objectives[pareto_indices]
        
        return pareto_front, pareto_objectives
    
    def _create_offspring(self, parents, bounds):
        """Create offspring through crossover and mutation."""
        offspring = []
        n_vars = len(bounds)
        
        for i in range(0, len(parents), 2):
            parent1 = parents[i]
            parent2 = parents[min(i + 1, len(parents) - 1)]
            
            # Crossover
            if np.random.random() < self.crossover_rate:
                # Simulated binary crossover (SBX)
                child1, child2 = self._sbx_crossover(
                    parent1, parent2, bounds
                )
            else:
                child1, child2 = parent1.copy(), parent2.copy()
            
            # Mutation
            if np.random.random() < self.mutation_rate:
                child1 = self._polynomial_mutation(child1, bounds)
            if np.random.random() < self.mutation_rate:
                child2 = self._polynomial_mutation(child2, bounds)
            
            offspring.extend([child1, child2])
        
        return offspring[:len(parents)]
    
    def _sbx_crossover(self, parent1, parent2, bounds, eta=20):
        """Simulated binary crossover."""
        child1 = parent1.copy()
        child2 = parent2.copy()
        
        for i in range(len(parent1)):
            if np.random.random() < 0.5:
                continue
            
            y1 = min(parent1[i], parent2[i])
            y2 = max(parent1[i], parent2[i])
            
            if abs(y2 - y1) < 1e-10:
                continue
            
            beta = 1.0 + (2.0 * (y1 - bounds[i][0]) / (y2 - y1))
            alpha = 2.0 - beta ** (-(eta + 1.0))
            
            u = np.random.random()
            if u <= 1.0 / alpha:
                beta_q = (u * alpha) ** (1.0 / (eta + 1.0))
            else:
                beta_q = (1.0 / (2.0 - u * alpha)) ** \
                        (1.0 / (eta + 1.0))
            
            child1[i] = 0.5 * ((y1 + y2) - beta_q * (y2 - y1))
            child2[i] = 0.5 * ((y1 + y2) + beta_q * (y2 - y1))
            
            # Ensure bounds
            child1[i] = np.clip(child1[i], bounds[i][0], bounds[i][1])
            child2[i] = np.clip(child2[i], bounds[i][0], bounds[i][1])
        
        return child1, child2
    
    def _polynomial_mutation(self, individual, bounds, eta=20):
        """Polynomial mutation."""
        mutated = individual.copy()
        
        for i in range(len(individual)):
            if np.random.random() > self.mutation_rate:
                continue
            
            y = individual[i]
            delta1 = (y - bounds[i][0]) / (bounds[i][1] - bounds[i][0])
            delta2 = (bounds[i][1] - y) / (bounds[i][1] - bounds[i][0])
            
            u = np.random.random()
            if u < 0.5:
                delta_q = (2.0 * u + (1.0 - 2.0 * u) * 
                          (1.0 - delta1) ** (eta + 1.0)) ** \
                          (1.0 / (eta + 1.0)) - 1.0
            else:
                delta_q = 1.0 - (2.0 * (1.0 - u) + 2.0 * 
                          (u - 0.5) * (1.0 - delta2) ** \
                          (eta + 1.0)) ** (1.0 / (eta + 1.0))
            
            mutated[i] = y + delta_q * (bounds[i][1] - bounds[i][0])
            mutated[i] = np.clip(mutated[i], bounds[i][0], bounds[i][1])
        
        return mutated
    
    def _select_next_generation(self, population, objectives):
        """Select next generation using elitism."""
        # Non-dominated sorting
        fronts = self.fast_non_dominated_sort(objectives)
        
        # Select individuals for next generation
        next_population = []
        next_objectives = []
        
        for front in fronts:
            if len(next_population) + len(front) <= self.population_size:
                # Add entire front
                next_population.extend([population[i] for i in front])
                next_objectives.extend([objectives[i] for i in front])
            else:
                # Add subset based on crowding distance
                n_remaining = self.population_size - len(next_population)
                distances = self.calculate_crowding_distance(
                    objectives, front
                )
                sorted_indices = np.argsort(distances)[::-1]
                selected = [front[i] for i in sorted_indices[:n_remaining]]
                
                next_population.extend([population[i] for i in selected])
                next_objectives.extend([objectives[i] for i in selected])
                break
        
        return next_population, np.array(next_objectives)
\end{lstlisting}

\section{UKF Sigma Point Generation}
\label{sec:ukf_sigma}

The Unscented Kalman Filter uses sigma points to capture the mean and covariance of the state distribution through nonlinear transformations.

\begin{lstlisting}[language=Python, caption={UKF Sigma Point Generation and Propagation}, label={lst:ukf_sigma}, numbers=left, numberstyle=\tiny, stepnumber=1, basicstyle=\ttfamily\footnotesize, breaklines=true, frame=single]
import numpy as np
from scipy.linalg import sqrtm, cholesky

class UnscentedKalmanFilter:
    """
    Unscented Kalman Filter for nonlinear state estimation
    of mobile robot pose and velocity.
    """
    def __init__(self, dim_x, dim_z, dt, alpha=1e-3, 
                 beta=2.0, kappa=0.0):
        """
        Initialize UKF.
        
        Args:
            dim_x: State dimension
            dim_z: Measurement dimension
            dt: Time step
            alpha: Sigma point spread parameter (1e-4 to 1)
            beta: Prior knowledge parameter (2 is optimal for Gaussian)
            kappa: Secondary scaling parameter (0 or 3-dim_x)
        """
        self.dim_x = dim_x
        self.dim_z = dim_z
        self.dt = dt
        
        # Sigma point parameters
        self.alpha = alpha
        self.beta = beta
        self.kappa = kappa
        
        # Compute lambda and weights
        self.lambda_ = alpha**2 * (dim_x + kappa) - dim_x
        
        # Number of sigma points
        self.n_sigma = 2 * dim_x + 1
        
        # Compute weights
        self.weights_m = np.zeros(self.n_sigma)
        self.weights_c = np.zeros(self.n_sigma)
        
        self.weights_m[0] = self.lambda_ / (dim_x + self.lambda_)
        self.weights_c[0] = self.weights_m[0] + (1 - alpha**2 + beta)
        
        for i in range(1, self.n_sigma):
            self.weights_m[i] = 0.5 / (dim_x + self.lambda_)
            self.weights_c[i] = self.weights_m[i]
        
        # State and covariance
        self.x = np.zeros(dim_x)  # State estimate
        self.P = np.eye(dim_x)    # Covariance matrix
        
        # Process and measurement noise covariances
        self.Q = np.eye(dim_x) * 0.01  # Process noise
        self.R = np.eye(dim_z) * 0.1   # Measurement noise
        
    def generate_sigma_points(self, x, P):
        """
        Generate sigma points around mean x with covariance P.
        
        Args:
            x: Mean state vector [dim_x]
            P: Covariance matrix [dim_x, dim_x]
            
        Returns:
            sigma_points: Array of sigma points [n_sigma, dim_x]
        """
        sigma_points = np.zeros((self.n_sigma, self.dim_x))
        
        # First sigma point is the mean
        sigma_points[0] = x
        
        # Compute matrix square root of scaled covariance
        # (dim_x + lambda) * P
        U = cholesky((self.dim_x + self.lambda_) * P).T
        
        # Generate remaining sigma points
        for i in range(self.dim_x):
            sigma_points[i + 1] = x + U[i]
            sigma_points[i + 1 + self.dim_x] = x - U[i]
        
        return sigma_points
    
    def unscented_transform(self, sigma_points, weights_m, 
                          weights_c, noise_cov=None):
        """
        Compute mean and covariance from transformed sigma points.
        
        Args:
            sigma_points: Transformed sigma points [n_sigma, dim]
            weights_m: Weights for mean computation
            weights_c: Weights for covariance computation
            noise_cov: Optional noise covariance to add
            
        Returns:
            mean: Weighted mean of sigma points
            cov: Weighted covariance of sigma points
        """
        # Compute weighted mean
        mean = np.dot(weights_m, sigma_points)
        
        # Compute weighted covariance
        deviations = sigma_points - mean
        cov = np.dot(weights_c * deviations.T, deviations)
        
        # Add noise covariance if provided
        if noise_cov is not None:
            cov += noise_cov
        
        return mean, cov
    
    def predict(self, control_input=None):
        """
        Prediction step of UKF.
        
        Args:
            control_input: Control vector [v, omega]
        """
        # Generate sigma points from current state
        sigma_points = self.generate_sigma_points(self.x, self.P)
        
        # Propagate sigma points through process model
        sigma_points_pred = np.zeros_like(sigma_points)
        for i in range(self.n_sigma):
            sigma_points_pred[i] = self.process_model(
                sigma_points[i], 
                control_input
            )
        
        # Compute predicted mean and covariance
        self.x, self.P = self.unscented_transform(
            sigma_points_pred,
            self.weights_m,
            self.weights_c,
            self.Q
        )
        
        # Store sigma points for update step
        self.sigma_points_pred = sigma_points_pred
    
    def update(self, measurement):
        """
        Update step of UKF.
        
        Args:
            measurement: Measurement vector [z1, z2, ...]
        """
        # Transform sigma points to measurement space
        sigma_points_meas = np.zeros((self.n_sigma, self.dim_z))
        for i in range(self.n_sigma):
            sigma_points_meas[i] = self.measurement_model(
                self.sigma_points_pred[i]
            )
        
        # Compute predicted measurement mean and covariance
        z_pred, S = self.unscented_transform(
            sigma_points_meas,
            self.weights_m,
            self.weights_c,
            self.R
        )
        
        # Compute cross-covariance
        Pxz = np.zeros((self.dim_x, self.dim_z))
        for i in range(self.n_sigma):
            dx = self.sigma_points_pred[i] - self.x
            dz = sigma_points_meas[i] - z_pred
            Pxz += self.weights_c[i] * np.outer(dx, dz)
        
        # Compute Kalman gain
        K = np.dot(Pxz, np.linalg.inv(S))
        
        # Update state estimate
        innovation = measurement - z_pred
        self.x = self.x + np.dot(K, innovation)
        
        # Update covariance estimate
        self.P = self.P - np.dot(K, np.dot(S, K.T))
        
        return innovation
    
    def process_model(self, x, u):
        """
        Nonlinear process model for mobile robot.
        
        State: [x, y, theta, v, omega]
        Control: [v_cmd, omega_cmd]
        
        Args:
            x: State vector
            u: Control input
            
        Returns:
            x_next: Next state
        """
        if u is None:
            u = np.zeros(2)
        
        # Extract state variables
        pos_x, pos_y, theta, v, omega = x
        v_cmd, omega_cmd = u
        
        # Kinematic model (with velocity dynamics)
        # Position update
        x_next = np.zeros_like(x)
        x_next[0] = pos_x + v * np.cos(theta) * self.dt
        x_next[1] = pos_y + v * np.sin(theta) * self.dt
        x_next[2] = theta + omega * self.dt
        
        # Velocity dynamics (first-order lag)
        tau_v = 0.2  # Time constant for linear velocity
        tau_omega = 0.15  # Time constant for angular velocity
        
        x_next[3] = v + (v_cmd - v) / tau_v * self.dt
        x_next[4] = omega + (omega_cmd - omega) / tau_omega * self.dt
        
        # Normalize angle to [-pi, pi]
        x_next[2] = np.arctan2(np.sin(x_next[2]), np.cos(x_next[2]))
        
        return x_next
    
    def measurement_model(self, x):
        """
        Measurement model (observes position and orientation).
        
        Args:
            x: State vector
            
        Returns:
            z: Measurement vector [x, y, theta]
        """
        # Direct observation of position and orientation
        z = x[:3]  # [x, y, theta]
        return z
    
    def set_process_noise(self, Q):
        """Set process noise covariance matrix."""
        self.Q = Q
    
    def set_measurement_noise(self, R):
        """Set measurement noise covariance matrix."""
        self.R = R
    
    def get_state(self):
        """Get current state estimate."""
        return self.x.copy()
    
    def get_covariance(self):
        """Get current covariance estimate."""
        return self.P.copy()
    
    def reset(self, x0, P0):
        """
        Reset filter with new initial state and covariance.
        
        Args:
            x0: Initial state
            P0: Initial covariance
        """
        self.x = x0.copy()
        self.P = P0.copy()

# Example usage
if __name__ == "__main__":
    # Initialize UKF for 5D state and 3D measurement
    dim_x = 5  # [x, y, theta, v, omega]
    dim_z = 3  # [x, y, theta]
    dt = 0.1
    
    ukf = UnscentedKalmanFilter(dim_x, dim_z, dt)
    
    # Set initial state
    x0 = np.array([0.0, 0.0, 0.0, 0.0, 0.0])
    P0 = np.eye(dim_x) * 0.1
    ukf.reset(x0, P0)
    
    # Set noise covariances
    Q = np.diag([0.01, 0.01, 0.01, 0.05, 0.05])
    R = np.diag([0.1, 0.1, 0.05])
    ukf.set_process_noise(Q)
    ukf.set_measurement_noise(R)
    
    # Simulate prediction and update
    control = np.array([0.5, 0.1])  # [v_cmd, omega_cmd]
    measurement = np.array([0.05, 0.005, 0.01])  # [x, y, theta]
    
    ukf.predict(control)
    ukf.update(measurement)
    
    print("State estimate:", ukf.get_state())
    print("Covariance:", ukf.get_covariance())
\end{lstlisting}

\section{DWA Trajectory Scoring}
\label{sec:dwa_scoring}

The Dynamic Window Approach generates and scores candidate trajectories for local obstacle avoidance and control.

\begin{lstlisting}[language=Python, caption={DWA Trajectory Generation and Scoring}, label={lst:dwa_scoring}, numbers=left, numberstyle=\tiny, stepnumber=1, basicstyle=\ttfamily\footnotesize, breaklines=true, frame=single]
import numpy as np
from typing import Tuple, List

class DynamicWindowApproach:
    """
    Dynamic Window Approach for local trajectory planning.
    Generates and evaluates trajectories within robot's 
    dynamic constraints.
    """
    def __init__(self, config=None):
        """
        Initialize DWA planner.
        
        Args:
            config: Configuration dictionary with DWA parameters
        """
        if config is None:
            config = {}
        
        # Robot constraints
        self.max_speed = config.get('max_speed', 1.0)  # m/s
        self.min_speed = config.get('min_speed', 0.0)
        self.max_yaw_rate = config.get('max_yaw_rate', 1.0)  # rad/s
        self.max_accel = config.get('max_accel', 0.5)  # m/s^2
        self.max_yaw_accel = config.get('max_yaw_accel', 1.0)  # rad/s^2
        
        # Trajectory evaluation parameters
        self.dt = config.get('dt', 0.1)  # Time step
        self.predict_time = config.get('predict_time', 3.0)  # Prediction horizon
        self.v_resolution = config.get('v_resolution', 0.1)  # Velocity sampling
        self.yaw_rate_resolution = config.get('yaw_rate_resolution', 0.1)
        
        # Scoring weights
        self.heading_weight = config.get('heading_weight', 1.0)
        self.distance_weight = config.get('distance_weight', 1.0)
        self.velocity_weight = config.get('velocity_weight', 1.0)
        self.energy_weight = config.get('energy_weight', 0.5)
        
        # Safety parameters
        self.robot_radius = config.get('robot_radius', 0.3)  # m
        self.collision_threshold = config.get('collision_threshold', 0.5)
    
    def plan(self, state, goal, obstacles):
        """
        Plan optimal trajectory using DWA.
        
        Args:
            state: Current robot state [x, y, theta, v, omega]
            goal: Goal position [x, y]
            obstacles: List of obstacles [[x, y, r], ...]
            
        Returns:
            best_u: Best control [v, omega]
            best_trajectory: Best predicted trajectory
        """
        # Generate dynamic window
        dw = self.calculate_dynamic_window(state)
        
        # Generate all candidate trajectories
        trajectories, controls = self.generate_trajectories(state, dw)
        
        # Score all trajectories
        scores = self.score_trajectories(
            trajectories, 
            controls,
            goal, 
            obstacles
        )
        
        # Select best trajectory
        if len(scores) == 0:
            # Emergency stop if no valid trajectory
            return np.array([0.0, 0.0]), np.array([state[:3]])
        
        best_idx = np.argmax(scores)
        best_u = controls[best_idx]
        best_trajectory = trajectories[best_idx]
        
        return best_u, best_trajectory
    
    def calculate_dynamic_window(self, state):
        """
        Calculate dynamic window based on current velocity and 
        robot dynamics.
        
        Args:
            state: Current state [x, y, theta, v, omega]
            
        Returns:
            dw: Dynamic window [v_min, v_max, omega_min, omega_max]
        """
        x, y, theta, v, omega = state
        
        # Velocity limits based on robot constraints
        v_min = self.min_speed
        v_max = self.max_speed
        omega_min = -self.max_yaw_rate
        omega_max = self.max_yaw_rate
        
        # Dynamic window based on acceleration limits
        v_min_dyn = v - self.max_accel * self.dt
        v_max_dyn = v + self.max_accel * self.dt
        omega_min_dyn = omega - self.max_yaw_accel * self.dt
        omega_max_dyn = omega + self.max_yaw_accel * self.dt
        
        # Intersect robot limits with dynamic limits
        v_min = max(v_min, v_min_dyn)
        v_max = min(v_max, v_max_dyn)
        omega_min = max(omega_min, omega_min_dyn)
        omega_max = min(omega_max, omega_max_dyn)
        
        dw = [v_min, v_max, omega_min, omega_max]
        return dw
    
    def generate_trajectories(self, state, dw):
        """
        Generate candidate trajectories within dynamic window.
        
        Args:
            state: Current state [x, y, theta, v, omega]
            dw: Dynamic window [v_min, v_max, omega_min, omega_max]
            
        Returns:
            trajectories: List of trajectories
            controls: List of control inputs
        """
        trajectories = []
        controls = []
        
        # Sample velocities within dynamic window
        v_samples = np.arange(dw[0], dw[1], self.v_resolution)
        omega_samples = np.arange(dw[2], dw[3], self.yaw_rate_resolution)
        
        # Generate trajectory for each control sample
        for v in v_samples:
            for omega in omega_samples:
                # Simulate trajectory
                trajectory = self.simulate_trajectory(
                    state, 
                    np.array([v, omega])
                )
                trajectories.append(trajectory)
                controls.append(np.array([v, omega]))
        
        return trajectories, controls
    
    def simulate_trajectory(self, state, control):
        """
        Simulate trajectory for given control input.
        
        Args:
            state: Initial state [x, y, theta, v, omega]
            control: Control input [v, omega]
            
        Returns:
            trajectory: Predicted trajectory points
        """
        x, y, theta, v, omega = state
        v_cmd, omega_cmd = control
        
        trajectory = []
        n_steps = int(self.predict_time / self.dt)
        
        for _ in range(n_steps):
            # Update state with simple kinematic model
            x += v * np.cos(theta) * self.dt
            y += v * np.sin(theta) * self.dt
            theta += omega * self.dt
            
            # First-order velocity dynamics
            v += (v_cmd - v) * 0.5 * self.dt
            omega += (omega_cmd - omega) * 0.5 * self.dt
            
            trajectory.append([x, y, theta])
        
        return np.array(trajectory)
    
    def score_trajectories(self, trajectories, controls, 
                          goal, obstacles):
        """
        Score all candidate trajectories.
        
        Args:
            trajectories: List of trajectory arrays
            controls: List of control inputs
            goal: Goal position [x, y]
            obstacles: List of obstacles [[x, y, r], ...]
            
        Returns:
            scores: Score for each trajectory
        """
        scores = []
        
        for trajectory, control in zip(trajectories, controls):
            # Check collision
            if self.check_collision(trajectory, obstacles):
                continue  # Skip colliding trajectories
            
            # Compute individual scores
            heading_score = self.heading_cost(trajectory, goal)
            distance_score = self.distance_cost(trajectory, goal)
            velocity_score = self.velocity_cost(control)
            energy_score = self.energy_cost(control)
            
            # Combine scores with weights
            total_score = (
                self.heading_weight * heading_score +
                self.distance_weight * distance_score +
                self.velocity_weight * velocity_score -
                self.energy_weight * energy_score
            )
            
            scores.append(total_score)
        
        return np.array(scores)
    
    def heading_cost(self, trajectory, goal):
        """
        Score based on heading toward goal.
        
        Args:
            trajectory: Trajectory points
            goal: Goal position [x, y]
            
        Returns:
            score: Heading score (higher is better)
        """
        # Use final point in trajectory
        final_point = trajectory[-1]
        x, y, theta = final_point
        
        # Angle to goal
        angle_to_goal = np.arctan2(goal[1] - y, goal[0] - x)
        
        # Angular difference
        angle_diff = np.abs(self.normalize_angle(angle_to_goal - theta))
        
        # Convert to score (0 to 1, where 1 is perfect alignment)
        score = 1.0 - (angle_diff / np.pi)
        
        return score
    
    def distance_cost(self, trajectory, goal):
        """
        Score based on distance to goal.
        
        Args:
            trajectory: Trajectory points
            goal: Goal position [x, y]
            
        Returns:
            score: Distance score (higher is better)
        """
        # Use final point in trajectory
        final_point = trajectory[-1]
        x, y = final_point[:2]
        
        # Distance to goal
        distance = np.hypot(goal[0] - x, goal[1] - y)
        
        # Convert to score (normalize and invert)
        # Closer is better
        max_distance = 10.0  # Maximum expected distance
        score = 1.0 - min(distance / max_distance, 1.0)
        
        return score
    
    def velocity_cost(self, control):
        """
        Score based on velocity (prefer higher velocities).
        
        Args:
            control: Control input [v, omega]
            
        Returns:
            score: Velocity score (higher is better)
        """
        v, omega = control
        
        # Normalize to [0, 1]
        score = v / self.max_speed
        
        return score
    
    def energy_cost(self, control):
        """
        Estimate energy cost for control input.
        
        Args:
            control: Control input [v, omega]
            
        Returns:
            cost: Energy cost (lower is better)
        """
        v, omega = control
        
        # Simplified energy model
        # Energy proportional to v^2 + omega^2
        k_linear = 0.5
        k_angular = 0.3
        
        energy = k_linear * v**2 + k_angular * omega**2
        
        # Normalize to [0, 1]
        max_energy = (k_linear * self.max_speed**2 + 
                     k_angular * self.max_yaw_rate**2)
        cost = energy / max_energy
        
        return cost
    
    def check_collision(self, trajectory, obstacles):
        """
        Check if trajectory collides with any obstacle.
        
        Args:
            trajectory: Trajectory points
            obstacles: List of obstacles [[x, y, r], ...]
            
        Returns:
            collision: True if collision detected
        """
        for point in trajectory:
            x, y = point[:2]
            
            for obstacle in obstacles:
                obs_x, obs_y, obs_r = obstacle
                distance = np.hypot(x - obs_x, y - obs_y)
                
                # Check if within collision threshold
                if distance < (self.robot_radius + obs_r + 
                              self.collision_threshold):
                    return True
        
        return False
    
    def normalize_angle(self, angle):
        """Normalize angle to [-pi, pi]."""
        while angle > np.pi:
            angle -= 2.0 * np.pi
        while angle < -np.pi:
            angle += 2.0 * np.pi
        return angle


# Example usage
if __name__ == "__main__":
    # Configure DWA
    config = {
        'max_speed': 1.0,
        'max_yaw_rate': 1.0,
        'max_accel': 0.5,
        'max_yaw_accel': 1.0,
        'v_resolution': 0.1,
        'yaw_rate_resolution': 0.1,
        'dt': 0.1,
        'predict_time': 3.0,
        'heading_weight': 1.0,
        'distance_weight': 1.0,
        'velocity_weight': 0.5,
        'energy_weight': 0.3,
        'robot_radius': 0.3
    }
    
    dwa = DynamicWindowApproach(config)
    
    # Test planning
    state = np.array([0.0, 0.0, 0.0, 0.0, 0.0])
    goal = np.array([5.0, 5.0])
    obstacles = [
        [2.0, 2.0, 0.5],
        [3.0, 1.0, 0.3]
    ]
    
    control, trajectory = dwa.plan(state, goal, obstacles)
    print("Optimal control:", control)
    print("Trajectory shape:", trajectory.shape)
\end{lstlisting}

\clearpage
